\section{Introduction}
\label{sec:introduction}

Magnetic micro/nanomanipulation using electromagnetic actuators has become
a powerful tool for biological and material science research. By generating
spatially controlled magnetic field gradients, these actuators exert piconewton-
to nanonewton-scale forces on paramagnetic microbeads, enabling the study of
cellular mechanics, molecular bonds, and surface interactions.

The design of such actuators has been advanced significantly by
Zhang and Menq~\cite{Zhang2010}, who developed a four-pole (quadrupole)
magnetic force actuator with a ferromagnetic yoke, and later by
Menq~\cite{Menq2011}, who extended the design to a six-pole (hexapole)
configuration for full three-dimensional force control. In both designs,
the ferromagnetic yoke serves as a low-reluctance return path for the
magnetic flux, concentrating the field at the pole tips and enabling
strong field gradients over short distances.

A central parameter in these designs is the \emph{magnetic charge offset}
$b$, which characterizes how far behind the pole tip surface the equivalent
magnetic monopole charge resides. Physically, $b$ encapsulates the combined
effect of pole geometry, magnetic circuit reluctance, and flux concentration.
A smaller $b$ implies tighter flux concentration and stronger gradients at a
given distance, making $b$ the appropriate figure of merit for comparing
different actuator geometries.

In this work, we characterize a yokeless hexapole actuator---a configuration
where the six poles operate without a ferromagnetic yoke, so that the flux
return path passes entirely through air. This fundamentally different magnetic
circuit topology is expected to yield a larger $b$ (weaker flux concentration)
compared to yoke-based designs. The primary objective is to determine $b$
experimentally through frequency-domain measurements and to compare the
result with the yoke-based quadrupole of Zhang and Menq~\cite{Zhang2010},
for which $b \approx \SI{405}{\um}$.

The remainder of this report is organized as follows.
\Cref{sec:theory} develops the monopole model and derives the
signal-chain transfer function $H(x)$.
\Cref{sec:experiment} describes the experimental setup and measurement protocol.
\Cref{sec:results} presents the fitting results and discusses frequency
independence and distance range sensitivity.
\Cref{sec:physics} derives the physical parameters ($\kpole$, $\Ra$) and
estimates the magnetic force on a paramagnetic bead.
\Cref{sec:discussion} compares the yokeless and yoke-based designs.
\Cref{sec:conclusion} summarizes the findings.
